\documentclass[10pt,a4paper]{amsart}
\usepackage[margin=19mm]{geometry}
\usepackage{amsmath,amssymb,mathtools}
\usepackage[scr=rsfs,cal=euler]{mathalfa}
\usepackage{xfrac}
\usepackage[unicode,hidelinks,bookmarks=false]{hyperref}
\newcommand{\ie}{i.e.\ }
\newcommand{\cf}{cf.\ }
\newcommand{\resp}{respectively\ }
\newcommand{\Subsection}{\subsection}
\providecommand{\nobreakdash}{\nobreak\hbox{-}}
\setlength{\emergencystretch}{2em}
\multlinegap0pt
\usepackage{bm}
\usepackage{bbm}
\usepackage{dsfont}
\usepackage{xspace}
\usepackage{cancel}
\usepackage{mathtools}
\usepackage{amsthm}
\makeatletter
\@ifundefined{theorem}{\newtheorem{theorem}{Theorem}}{}
\@ifundefined{definition}{\newtheorem{definition}[theorem]{Definition}}{}
\@ifundefined{lemma}{\newtheorem{lemma}[theorem]{Lemma}}{}
\makeatother
\def\N{\mathbb{N}}
\def\R{\mathcal{R}}
\def\S{\mathfrak{S}}
\def\Z{\mathbb{Z}}
\def\dmax{\kappa}
\def\hexpl#1{\hbox{ $\langle$\textit{#1}$\rangle$}}
\def\zmax{\nu}
\def\zmin{\mu}
\title[Two Useful Facts About Generating Functions]{Two Useful Facts About Generating Functions}
\author{Alex Kasman, Robert Milson}
\date{Source: arXiv:2511.11559v1 (CC BY 4.0)}
\begin{document}
\maketitle

\noindent\emph{Source and licence.} Two Useful Facts About Generating Functions. Version arXiv:2511.11559v1 (CC BY 4.0), distributed under the Creative Commons Attribution 4.0 International licence, as recorded in the arXiv licence metadata for this exact version and shown on the abstract page https://arxiv.org/abs/2511.11559v1. Full rights notice: licenses/arxiv-2511.11559-CC-BY-4.0.txt.\par\smallskip

\noindent\textit{Editorial scope.} Counted unit: the Lemma whose source label is \texttt{lem:main} in the article (source0.tex lines 278--284, source label \texttt{lem:main}), delivered together with the article's own complete original proof of it (source0.tex lines 285--340, one \texttt{proof} environment of 56 lines). No mathematics is written, repaired or re-derived: statement and proof are the article's own text line for line. Required context retained uncounted: def:LOmega (lines 250--266). Every definition, display and label of those lines is retained unchanged. Omitted from this package and not redistributed: 1--249, 267--277, 341--931. Nothing inside a retained excerpt is omitted or altered: every retained body is delivered exactly as the article's lines are written, so no omission receipt of any kind accompanies this package. Citations: the retained excerpts carry no citation command at all, so no bibliography entry of the article is transcribed and no reference is redistributed. Reference adaptation: none was needed and none was made; every reference inside the delivered unit resolves inside it, so no unresolved reference or citation is produced. Printed slips kept literal: no printed word, symbol, hypothesis or number of the counted statement or of its proof was altered. Layout and class: the article's own class is \texttt{amsart} with packages including amssymb; the wrapper is the lane's pinned \texttt{amsart} editorial wrapper at 10pt on A4, which restates the theorem-like environments the retained text needs and adds the article's own macro definitions that the retained text needs (\texttt{\textbackslash N}, \texttt{\textbackslash R}, \texttt{\textbackslash S}, \texttt{\textbackslash Z}, \texttt{\textbackslash dmax}, \texttt{\textbackslash hexpl}, \texttt{\textbackslash zmax}, \texttt{\textbackslash zmin}), with extra wrapper packages (bm, bbm, dsfont, xspace, cancel, mathtools). The delivered numbering is the unit's own. Assets: no retained span contains a figure environment, a graphics-inclusion command, a PDF-page inclusion, a table or a TikZ picture; no image file is copied into this package, the package declares no asset dependency and no image file is redistributed with it. Licence: version arXiv:2511.11559v1 of this article is distributed under the Creative Commons Attribution 4.0 International licence, as recorded in the arXiv licence metadata for this exact version and shown on the abstract page https://arxiv.org/abs/2511.11559v1. A full rights notice is delivered at licenses/arxiv-2511.11559-CC-BY-4.0.txt. Characters: every retained excerpt is pure ASCII, so the delivered engine, XeLaTeX with its default Latin Modern text font, reports no missing character.
\begin{definition}\label{def:LOmega}
Suppose $L\in \R_0[z,z^{-1},\partial_z]$ is an ordinary differential
operator in $z$ with singularities only at $z=0$.   Then it can be written in the form
$$
L=\sum_{i=\zmin}^{\zmax}\sum_{j=0}^{\dmax} c_{ij}z^{i}\partial_z^{j},
\qquad
\zmin<\zmax\in\Z,\ \dmax\in\N\cup\{0\}, c_{ij}\in\R_0.
$$
Associate to the choice of any such differential operator a
corresponding \textit{difference} operator $\Omega_L$ defined as
\begin{eqnarray*}
\Omega_L&=&\sum_{i=\zmin}^{\zmax}\sum_{j=0}^{\dmax} c_{ij}
\prod_{k=0}^{j-1}(n+j-i-k) \S^{j-i}  \\
  &=&\sum_{i=\zmin}^{\zmax}\sum_{j=0}^{\dmax} c_{ij}
(n+j-i)(n+j-i-1)\cdots(n-i+1) \S^{j-i}.
\end{eqnarray*}
\end{definition}

\begin{lemma}\label{lem:main} For any differential operator $L\in\R_0[z,z^{-1},\partial_z]$ and
  the corresponding difference operator $\Omega_L$ as in
  Definition~\ref{def:LOmega} we have
$$
L(\psi)=\sum_{n=\zmin}^{\infty} \Omega_L(a)_nz^n.
$$
\end{lemma}

\begin{proof}
For $j,m\in\N\cup\{0\}$ one has
\begin{equation}\label{eqn:simple}
\partial_z^{j}(z^m)
=
m(m-1)\cdots(m-j+1)z^{m-j}.
\end{equation}
Note that this formula is valid even in the case that $m<j$ because in
that case $z^m$ is in the kernel of the operator and a factor of
$(m-m)$ appears in the product on the right side, so they are both zero.

It then follows by linearity that
\begin{eqnarray*}
  L(\psi)&=&\sum_{i=\zmin}^{\zmax}\sum_{j=0}^{\dmax}c_{ij}z^i\partial_z^j\psi(z)\\
         &=&\sum_{i=\zmin}^{\zmax}\sum_{j=0}^{\dmax}
             c_{ij}z^i\partial_z^j\left( \sum_{m=0}^\infty a_m z^m\right)
  \hexpl{expanding $\psi$ using the index $m$}\\
         &=&\sum_{i=\zmin}^{\zmax}\sum_{j=0}^{\dmax}\sum_{m=0}^\infty c_{ij}a_mz^i\left(\partial_z^jz^m\right)\\
  &=&\sum_{i=\zmin}^{\zmax}\sum_{j=0}^{\dmax}\sum_{m=0}^\infty c_{ij}a_mz^i\left( m(m-1)\cdots(m-j+1)z^{m-j}\right)
\hexpl{by \eqref{eqn:simple}}\\
%\end{eqnarray*}
%\begin{eqnarray*}
\phantom{L(\psi)}  &=&\sum_{i=\zmin}^{\zmax}\sum_{j=0}^{\dmax}\sum_{m=0}^\infty c_{ij}a_m
      m(m-1)\cdots(m-j+1)z^{m-j+i}\\
%&&\hexpl{combining powers of $z$}\\
  &=&\sum_{i=\zmin}^{\zmax}\sum_{j=0}^{\dmax}\sum_{n=i-j}^\infty c_{ij}a_{n+j-i}
      (n+j-i)(n+j-i-1)\cdots(n-i+1)z^{n}.
%&&\hexpl{rewriting in terms of $n=m-j+i$}.
\end{eqnarray*}

The final equality above is obtained by rewriting the inner-most sum
in terms of the new variable $n=m-j+i$.
After doing so, this sum begins at $n=i-j$.  However, the equality would
still be true if it started at a value \textit{lower} than $i-j$ since
$a_n=0$ for $n<0$.  In addition, the smallest power of $z$ which can
appear in this expression with a non-zero coefficient is $z^{\zmin}$
because $\psi$ is expanded only in non-negative powers of $z$, their
derivatives only involve non-negative powers of $z$, and those are
multiplied by $z^i$ for $i\geq \zmin$.  So, the sum in $n$ can be
taken to start at $n=\zmin$.  Then, since there are only a finite
number of terms in these sums having $z^n$ for any fixed value of $n$,
there can be no issues about reordering that would prevent us from
reordering the sums.
Consequently, we have
\begin{eqnarray*}
L(\psi)& =&\sum_{i=\zmin}^{\zmax}\sum_{j=0}^{\dmax}\sum_{n=\zmin}^\infty c_{ij}a_{n+j-i}
      (n+j-i)(n+j-i-1)\cdots(n-i+1)z^{n}\\
&=&\sum_{n=\zmin}^{\infty}\left(
\sum_{i=\zmin}^{\zmax}\sum_{j=0}^{\dmax} 
c_{ij}(n+j-i)(n+j-i-1)\cdots(n-i+1)a_{n+j-i}\right)z^{n}.\\
%&=&\sum_{n=\zmin}^{\infty}\Omega_L(a)_nz^n.
\end{eqnarray*}
The claim follows from the observation that the coefficient of each
term $z^n$
in the formula above is precisely $\Omega_L(a)_n$.
\end{proof}

\end{document}
